\documentclass[../main.tex]{subfiles}

\begin{document}

\section{Organizar las funciones/scripts de análsis de resultados}

Pues tengo un caos.

\subsection{Definición de protocolo}

El flujo general del script es el siguiente

\begin{enumerate}
    \item Enlistar todos los sistemas disponibles a analizar
    \item Seleccionar el sistema a analizar
    \item Extraer el archivo dat del sistema a analizar
    \item Analizar la información
\end{enumerate}

La selección del sistema es en base a los siguientes parámetros, densidad de empaquetamiento, temperatura, número de partículas y concentración de \emph{cross-linkers}.

\subsection{Metodología de los scripts}

En esta sección escribiré el flujo de información que voy formando.

\begin{enumerate}
    \item Obtener los directorios donde se almacenan los resultados numéricos 
    \item Extraer la información del archivo dat\footnote{Cuando se extrae se genera una nueva columna con el directorio del archivo} que hay cada uno de los directorios
    \item Agrupar por parámetros del sistema: \(\phi,T,\mathrm{N}_{\mathrm{part}},\%\mathrm{CL},\mathrm{damp},dt\).
    \item Crear un \emph{id} en términos de los parámetros anteriores
\end{enumerate}

\end{document}
